\chapter{Conclusion}
\label{ch:conclusion}

\section{What was built}

A C++23 numerical library in which twelve iterative solvers are written once,
against one problem interface, and run unchanged over seven execution backends.
Around them: a numerics module with root finders, quadrature rules, ODE
integrators and dense factorisations, each with its convergence order tested
against theory; two problem families; a benchmark harness that resolves a
declared matrix, resumes after interruption and records the provenance of every
row; and three reports built from the measurements by script.

\section{What was established}

\begin{enumerate}
  \item \textbf{Identical numerics across execution models is achievable, not
        merely approachable.} With deterministic reductions every shared memory
        backend at every worker count produces bit identical iterates, and the
        GPU sweeps match the CPU bit for bit as well. This is asserted by the
        test suite as exact equality across twelve solvers, two problem
        families, four backends and seven worker counts, and it required fixing
        the reduction chunk grid by problem size and disabling fused multiply
        add contraction on both host and device.

  \item \textbf{The classical convergence theory holds to the digit.} The Jacobi
        spectral radius matches $\cos(\pi h)$; Gauss Seidel needs exactly half
        the iterations of Jacobi, as $\rho_{GS} = \rho_J^2$ requires; Young's
        closed form relaxation factor is a genuine minimum; optimal relaxation
        changes the growth order from $O(n^2)$ to $O(n)$; and conjugate gradient
        reaches the same order without being told a parameter. The five point
        stencil is second order accurate with the error constant $\pi^2/12$
        predicted by the leading truncation term.

  \item \textbf{Sequential methods do not parallelise, and saying so is more
        useful than hiding it.} Natural ordering Gauss Seidel keeps exact
        semantics on every backend, including across MPI ranks, and shows
        essentially no speedup. Red black ordering is what buys the parallelism,
        at a measured cost of 2.6 percent more iterations.

  \item \textbf{A GPU speedup on a bandwidth bound kernel is mostly the memory
        system.} Decomposed through \eqref{eq:decomposition}, the observed
        advantage separates into a large factor that is the difference in
        measured bandwidth and a small factor that is the difference in how well
        each device is used. Reporting the undecomposed ratio would have been
        accurate and misleading.
\end{enumerate}

\section{What was not established}

Stated as plainly as the conclusions, because a limitation left implicit reads
as a claim.

\begin{itemize}
  \item \textbf{The performance versus efficiency core split could not be
        identified from inside the guest.} WSL2 synthesises a homogeneous
        topology and affinity binds to a virtual processor. The knee is reported
        from the aggregate scaling curve instead, which is a weaker claim than
        originally intended and the strongest the environment supports.
  \item \textbf{Nothing is established about compute bound problems.} Every
        conclusion about devices here concerns bandwidth bound stencil work. A
        dense factorisation would rank the hardware differently.
  \item \textbf{Nothing is established about single solve latency}, power, cost,
        or programmer effort. Transfer time is reported separately rather than
        folded in; the rest were not measured and are not claimed.
  \item \textbf{The iteration counts at the largest grids are extrapolated}, not
        measured, for the methods whose counts grow like $n^2$. The
        extrapolation rests on closed form rates that were verified at the sizes
        where measurement was affordable, and the report says which numbers are
        which.
\end{itemize}

\section{Further work}

\paragraph{Preconditioned conjugate gradient.} The symmetric splittings are
already implemented and symmetric Gauss Seidel and SSOR are exactly the
preconditioners the method wants. Combining them would be a small amount of code
and would demonstrate the point of the symmetric variants, which currently exist
in the zoo without the application that motivates them.

\paragraph{Multigrid.} Every component is present: smoothers in the relaxation
sweeps, an exact coarse solve in the LU routine, and the grid structure in the
Poisson problem. Multigrid is the method that makes the iteration count
independent of the grid entirely, and adding it would put the $O(n^2)$ and
$O(n)$ results of this report in their proper context as the two things it
improves on.

\paragraph{A local slab MPI decomposition.} Each rank currently allocates the
whole grid. The communication volume is already correct, so this would not change
any timing reported here, but it would remove the memory ceiling that keeps the
distributed runs below the sizes the device comparison uses.

\paragraph{Bare metal measurement of the core knee.} The one objective the
environment prevented. Running the same sweep outside a hypervisor would settle
whether the knee in the aggregate curve coincides with the performance core
count, which is currently inferred rather than confirmed.

\paragraph{Mixed precision.} The stencil sweeps are bandwidth bound, so halving
the width of the iterate would come close to halving the time. Whether the
resulting convergence penalty is worth it is exactly the kind of question this
library's separation of iteration count from per iteration cost is built to
answer.
